\section{Ownership, Returns and the Pace of Transformation}\label{sec:regimes}
\IfFileExists{part2/tables/regimes_macros.tex}{% Generated by code/verify_part2_regimes.py; do not edit. XBOUNDS placeholders: 2.5, 5, 9.
\providecommand{\regXbOne}{2.5}
\providecommand{\regXbGainOne}{0.24}
\providecommand{\regXbGainMyOne}{0.16}
\providecommand{\regXbMyOne}{1.92}
\providecommand{\regXbLossPEOne}{0.00}
\providecommand{\regXbTwo}{5}
\providecommand{\regXbGainTwo}{0.64}
\providecommand{\regXbGainMyTwo}{0.44}
\providecommand{\regXbMyTwo}{3.85}
\providecommand{\regXbLossPETwo}{0.51}
\providecommand{\regXbThree}{9}
\providecommand{\regXbGainThree}{1.48}
\providecommand{\regXbGainMyThree}{1.01}
\providecommand{\regXbMyThree}{6.92}
\providecommand{\regXbLossPEThree}{3.30}
\providecommand{\regXbBandLo}{6.85}
\providecommand{\regXbBandHi}{8.91}
\providecommand{\regVzeroMin}{5.9}
}{}% XBOUNDS placeholder values (generated by code/verify_part2_regimes.py)
% Generated by code/verify_part2_acquirer.py from notes/case_bending_spoons_f1_extract.md and data/examples/invoice_chain.json (version 2026-10-29c); do not edit.
\providecommand{\acqAcqHOneSix}{\$2{,}286 million}
\providecommand{\acqAcqShareHOneFive}{99.9\%}
\providecommand{\acqAcqShareHOneSix}{99.8\%}
\providecommand{\acqAcqShareTwentyFive}{100.0\%}
\providecommand{\acqAcqShareTwentyFiveTwo}{99.97\%}
\providecommand{\acqAcqShareTwentyFour}{81.7\%}
\providecommand{\acqAirEq}{\$2.25 billion}
\providecommand{\acqAirEv}{\$1.285 billion}
\providecommand{\acqAmzQOne}{11.5\%}
\providecommand{\acqAmzTwentyFive}{11.5\%}
\providecommand{\acqAmzTwentyFour}{12.8\%}
\providecommand{\acqAnnFactor}{0.51}
\providecommand{\acqAnnPtsAnnLo}{0.10}
\providecommand{\acqAnnPtsHi}{0.47}
\providecommand{\acqAnnPtsLo}{0.19}
\providecommand{\acqAoimQOne}{51.2\%}
\providecommand{\acqAoimTwentyFive}{47\%}
\providecommand{\acqAoimTwentyFiveX}{46.9\%}
\providecommand{\acqAoimTwentyFour}{44.6\%}
\providecommand{\acqAolB}{\$1.5 billion}
\providecommand{\acqAolDebt}{\$2.8 billion}
\providecommand{\acqAolEmp}{214}
\providecommand{\acqAolOpm}{53\%}
\providecommand{\acqAolWarn}{108}
\providecommand{\acqApps}{800{,}000}
\providecommand{\acqBaseTwentyFive}{5\%}
\providecommand{\acqBcCut}{198}
\providecommand{\acqBcPrice}{\$233 million}
\providecommand{\acqBcPs}{1.2}
\providecommand{\acqBcPx}{\$4.45}
\providecommand{\acqBcStaff}{600}
\providecommand{\acqBoundFac}{0.51}
\providecommand{\acqBrkOne}{\$600 million}
\providecommand{\acqBrkStake}{4\%}
\providecommand{\acqBrkTwo}{\$330 million}
\providecommand{\acqCapClose}{\$25.7 billion}
\providecommand{\acqCapOpen}{\$18.4 billion}
\providecommand{\acqCapexHOneSix}{\$4.0 million}
\providecommand{\acqCapexRevTwentyFive}{0.04\%}
\providecommand{\acqCapexRevTwentyFour}{0.09\%}
\providecommand{\acqCapexShareHOneFive}{0.05\%}
\providecommand{\acqCapexShareHOneSix}{0.2\%}
\providecommand{\acqCffShare}{148\%}
\providecommand{\acqCffTwentyFive}{\$1{,}935.9 million}
\providecommand{\acqCfiAbsTwentyFive}{\$1{,}851.9 million}
\providecommand{\acqCfiQOne}{274.1\%}
\providecommand{\acqCfiTwentyFive}{141.8\%}
\providecommand{\acqCfiTwentyFour}{132.6\%}
\providecommand{\acqChCouponExpHi}{113.7}
\providecommand{\acqChCouponExpLo}{86.7}
\providecommand{\acqChCouponHi}{112.8}
\providecommand{\acqChCouponLo}{86.0}
\providecommand{\acqChPaidHi}{113.1}
\providecommand{\acqChPaidLo}{86.6}
\providecommand{\acqConvHi}{114}
\providecommand{\acqConvLevHi}{103}
\providecommand{\acqConvLevLo}{79}
\providecommand{\acqConvLo}{87}
\providecommand{\acqCouponEur}{8.625}
\providecommand{\acqCouponMarg}{9.625}
\providecommand{\acqCov}{4.00}
\providecommand{\acqCrfHi}{10.8\%}
\providecommand{\acqCrfLo}{14.6\%}
\providecommand{\acqCumOutTwentyFive}{2.23}
\providecommand{\acqCumOutTwentyFour}{1.58}
\providecommand{\acqDayOne}{39.7\%}
\providecommand{\acqDeclHi}{9\%}
\providecommand{\acqDeclLo}{5\%}
\providecommand{\acqDepPrevQOne}{1.5}
\providecommand{\acqDepQOne}{\$2.01 billion}
\providecommand{\acqDepRatio}{10}
\providecommand{\acqDepTwentyThree}{\$194 million}
\providecommand{\acqDilution}{6.1\%}
\providecommand{\acqDoubleTwentyFive}{8.3}
\providecommand{\acqDoubleZeroTwentyFive}{10.9}
\providecommand{\acqEbM}{\$500 million}
\providecommand{\acqEbPx}{\$4.50}
\providecommand{\acqEqRaised}{\$549 million}
\providecommand{\acqEqShare}{99\%}
\providecommand{\acqEquityBs}{\$1.06 billion}
\providecommand{\acqFlagCash}{25\%}
\providecommand{\acqFlagCorr}{17}
\providecommand{\acqFlagCov}{26\%}
\providecommand{\acqFlagCovFast}{41\%}
\providecommand{\acqFlagG}{50\%}
\providecommand{\acqFlagPf}{10\%}
\providecommand{\acqFteTotal}{2{,}300}
\providecommand{\acqFullAnnAol}{\ensuremath{-}0.18}
\providecommand{\acqFullAnnEb}{\ensuremath{+}0.11}
\providecommand{\acqFullAnnVim}{\ensuremath{-}0.66}
\providecommand{\acqFullPtsHi}{2.2}
\providecommand{\acqFullPtsLo}{0.7}
\providecommand{\acqGeoEu}{21\%}
\providecommand{\acqGeoNa}{65\%}
\providecommand{\acqGiAol}{\ensuremath{-}5.1\%}
\providecommand{\acqGiEvernote}{\ensuremath{-}1.0\%}
\providecommand{\acqGiRemini}{\ensuremath{-}13.9\%}
\providecommand{\acqGiStreamyard}{\ensuremath{-}9.4\%}
\providecommand{\acqGiTwentyFive}{\ensuremath{-}5.1\%}
\providecommand{\acqGiTwentyFour}{\ensuremath{-}9.4\%}
\providecommand{\acqGinHi}{0.10}
\providecommand{\acqGinLo}{0.03}
\providecommand{\acqGmaxHi}{12\%}
\providecommand{\acqGmaxLo}{1\%}
\providecommand{\acqGmaxOrgHi}{20\%}
\providecommand{\acqGmaxOrgLo}{9\%}
\providecommand{\acqGridLevMax}{42\%}
\providecommand{\acqGridLevMin}{34\%}
\providecommand{\acqGridMax}{1.82}
\providecommand{\acqGridMfiveA}{1.82}
\providecommand{\acqGridMfourA}{1.46}
\providecommand{\acqGridMfourB}{1.11}
\providecommand{\acqGridMthreeA}{1.09}
\providecommand{\acqGridPfourA}{1.11}
\providecommand{\acqGridPfourB}{0.92}
\providecommand{\acqGrowthQOne}{132.2\%}
\providecommand{\acqGrowthTwentyFive}{94.7\%}
\providecommand{\acqGrowthTwentyFour}{73.4\%}
\providecommand{\acqHarvestMult}{15.2}
\providecommand{\acqHarvestNew}{\$2{,}110}
\providecommand{\acqHarvestOld}{\$130}
\providecommand{\acqHarvestRise}{1{,}523\%}
\providecommand{\acqHazFive}{0.147}
\providecommand{\acqHazMean}{0.125}
\providecommand{\acqHazTen}{0.127}
\providecommand{\acqHireRate}{0.036\%}
\providecommand{\acqHires}{286}
\providecommand{\acqHurL}{65\%}
\providecommand{\acqHurU}{25\%}
\providecommand{\acqIcCpool}{\$185.0 thousand}
\providecommand{\acqIcCrun}{\$329.0 thousand}
\providecommand{\acqIcCspec}{\$144.0 thousand}
\providecommand{\acqIcF}{\$2.81 million}
\providecommand{\acqIcFloorFour}{$-\$570$ thousand}
\providecommand{\acqIcFloorFourCertain}{$-\$458$ thousand}
\providecommand{\acqIcFloorInf}{\$141 thousand}
\providecommand{\acqIcFloorVone}{\$408 thousand}
\providecommand{\acqIcFreuse}{\$2.62 million}
\providecommand{\acqIcLimNoReb}{\$20.9 thousand}
\providecommand{\acqIcNbe}{7}
\providecommand{\acqIcNmin}{2.1}
\providecommand{\acqIcNminNoReb}{11.2}
\providecommand{\acqIcNpv}{$-\$3.78$ million}
\providecommand{\acqIcTh}{4}
\providecommand{\acqIcThbe}{12}
\providecommand{\acqIcUFbe}{0.84}
\providecommand{\acqIcVarpi}{0.82}
\providecommand{\acqIcVarpiBe}{0.73}
\providecommand{\acqIcW}{$-\$164.1$ thousand}
\providecommand{\acqIcWzero}{\$164.9 thousand}
\providecommand{\acqIntAbsQOne}{\$93.2 million}
\providecommand{\acqIntAbsTwentyFive}{\$142.6 million}
\providecommand{\acqIntPaidHOneSix}{\$161.2 million}
\providecommand{\acqIntPaidRatioHOneFive}{90\%}
\providecommand{\acqIntPaidRatioHOneSix}{80\%}
\providecommand{\acqIntPaidShQOne}{14.8\%}
\providecommand{\acqIntPaidShTwentyFive}{10.0\%}
\providecommand{\acqIntPaidShTwentyFour}{2.9\%}
\providecommand{\acqIntPaidShTwentyThree}{4.3\%}
\providecommand{\acqIntPaidTwentyFive}{\$130.9 million}
\providecommand{\acqIntQOne}{15.5\%}
\providecommand{\acqIntTwentyFive}{10.9\%}
\providecommand{\acqIntTwentyFour}{4.9\%}
\providecommand{\acqIntangShareTwentyFour}{18.3\%}
\providecommand{\acqIntangTwentyFour}{\$162.5 million}
\providecommand{\acqIpoCompany}{\$1.0 billion}
\providecommand{\acqIpoCompanyShares}{34.4 million}
\providecommand{\acqIpoNet}{\$1.10 billion}
\providecommand{\acqIpoNetEst}{\$933 million}
\providecommand{\acqIpoNetFull}{\$1{,}079 million}
\providecommand{\acqIpoNewP}{34.40 million}
\providecommand{\acqIpoPx}{\$29.00}
\providecommand{\acqIpoShares}{57.97 million}
\providecommand{\acqIpoSharesP}{57.97 million}
\providecommand{\acqIpoTotal}{\$1.68 billion}
\providecommand{\acqIrrL}{50.1\%}
\providecommand{\acqIrrPe}{19.9\%}
\providecommand{\acqIrrU}{22.3\%}
\providecommand{\acqKAol}{2.3}
\providecommand{\acqKEb}{1.7}
\providecommand{\acqKHi}{3.2}
\providecommand{\acqKLo}{1.7}
\providecommand{\acqKVim}{3.2}
\providecommand{\acqKcapHi}{2.1}
\providecommand{\acqKcapLo}{1.2}
\providecommand{\acqKomootMax}{85\%}
\providecommand{\acqKomootStaff}{150}
\providecommand{\acqKucoin}{\$13.1 billion}
\providecommand{\acqLevQOne}{2.19}
\providecommand{\acqLevQTwo}{2.4}
\providecommand{\acqLevTwentyFive}{2.24}
\providecommand{\acqLiabShare}{85\%}
\providecommand{\acqLifeHi}{19.5}
\providecommand{\acqLifeLo}{10.6}
\providecommand{\acqLogGrowthTwentyFive}{66.6\%}
\providecommand{\acqLtv}{85\%}
\providecommand{\acqMaiQOne}{4.5\%}
\providecommand{\acqMaiTwentyFive}{10.4\%}
\providecommand{\acqMaiTwentyFour}{14.1\%}
\providecommand{\acqMaxConcurrent}{20}
\providecommand{\acqMaxConcurrentExact}{20.7}
\providecommand{\acqMcfUQOne}{24.4\%}
\providecommand{\acqMcfUTwentyFive}{30.5\%}
\providecommand{\acqMcfUTwentyFour}{34.2\%}
\providecommand{\acqMgridHi}{0.3}
\providecommand{\acqMgridLo}{0.2}
\providecommand{\acqMiroEq}{\$1.79 billion}
\providecommand{\acqMiroEv}{\$1.355 billion}
\providecommand{\acqMiroReinv}{\$295 million}
\providecommand{\acqMixFast}{0.54}
\providecommand{\acqMixNrr}{83\%}
\providecommand{\acqMixSlow}{0.10}
\providecommand{\acqMixW}{22\%}
\providecommand{\acqMosaic}{330}
\providecommand{\acqMwBunch}{9.2\%}
\providecommand{\acqMwBunchLate}{31.4\%}
\providecommand{\acqMwCross}{2.46\%}
\providecommand{\acqMwCrossRise}{4\%}
\providecommand{\acqMwLate}{9.0\%}
\providecommand{\acqMwOne}{2.4\%}
\providecommand{\acqMwOneExact}{2.37\%}
\providecommand{\acqMwSpread}{9.5\%}
\providecommand{\acqMwZ}{0.176}
\providecommand{\acqMwZdays}{3}
\providecommand{\acqNeedTwentyFive}{9\%}
\providecommand{\acqNeedTwentyFour}{11\%}
\providecommand{\acqNeededTwentyFive}{99.7\%}
\providecommand{\acqNetDebt}{\$4.09 billion}
\providecommand{\acqNiTwentyFive}{\$\ensuremath{-}0.2 million}
\providecommand{\acqNotesCoupon}{168}
\providecommand{\acqNotesDue}{2032}
\providecommand{\acqNotesEur}{EUR 500 million}
\providecommand{\acqNotesEurC}{8.625\%}
\providecommand{\acqNotesUsd}{\$1.25 billion}
\providecommand{\acqNotesUsdC}{9.625\%}
\providecommand{\acqNrrAol}{95\%}
\providecommand{\acqNrrEvernote}{99\%}
\providecommand{\acqNrrHi}{95\%}
\providecommand{\acqNrrLo}{91\%}
\providecommand{\acqNrrQOne}{94\%}
\providecommand{\acqNrrRemini}{87\%}
\providecommand{\acqNrrStreamyard}{91\%}
\providecommand{\acqNrrTwentyFive}{95\%}
\providecommand{\acqNrrTwentyFour}{91\%}
\providecommand{\acqNrrTwentyThree}{93\%}
\providecommand{\acqNuTwentyFive}{0.051}
\providecommand{\acqNuTwentyFour}{0.094}
\providecommand{\acqOcfCagr}{122\%}
\providecommand{\acqOcfCfiQOne}{\ensuremath{-}261.5\%}
\providecommand{\acqOcfCfiTwentyFive}{\ensuremath{-}119.5\%}
\providecommand{\acqOcfCfiTwentyFour}{\ensuremath{-}102.0\%}
\providecommand{\acqOcfGrowth}{246\%, 42\% and 254\%}
\providecommand{\acqOcfTwentyFive}{\$290.6 million}
\providecommand{\acqOcfmQOne}{12.6\%}
\providecommand{\acqOcfmTwentyFive}{22.2\%}
\providecommand{\acqOcfmTwentyFour}{30.5\%}
\providecommand{\acqOpmQOne}{20.0\%}
\providecommand{\acqOpmTwentyFive}{21.3\%}
\providecommand{\acqOpmTwentyFour}{19.0\%}
\providecommand{\acqOrgGTwentyFive}{13\%}
\providecommand{\acqOrgGTwentyFour}{7\%}
\providecommand{\acqOrgLogQTwo}{3.0\%}
\providecommand{\acqOrgPage}{80}
\providecommand{\acqOrgQOne}{83\%}
\providecommand{\acqOrgTwentyFive}{79\%}
\providecommand{\acqOrgTwentyFour}{76\%}
\providecommand{\acqOutPrevTwentyFive}{2.8}
\providecommand{\acqOutPrevTwentyFour}{2.3}
\providecommand{\acqPfExcess}{100\%}
\providecommand{\acqPfInt}{14.1\%}
\providecommand{\acqPfOpm}{19.5\%}
\providecommand{\acqPfPage}{143}
\providecommand{\acqPfRatio}{2.0}
\providecommand{\acqPfRev}{\$2{,}608.5 million}
\providecommand{\acqPninetyfiveHi}{13}
\providecommand{\acqPninetyfiveLo}{7}
\providecommand{\acqPrHi}{90\%}
\providecommand{\acqPrLo}{10\%}
\providecommand{\acqPrSolo}{70\%}
\providecommand{\acqPrivRound}{\$11 billion}
\providecommand{\acqPxAol}{\$1{,}454 million}
\providecommand{\acqPxEb}{\$505 million}
\providecommand{\acqPxVim}{\$1{,}359 million}
\providecommand{\acqQTwoAdjM}{54\%}
\providecommand{\acqQTwoAdjMx}{54.1\%}
\providecommand{\acqQTwoAmz}{11.7\%}
\providecommand{\acqQTwoG}{126\%}
\providecommand{\acqQTwoGx}{126.3\%}
\providecommand{\acqQTwoInt}{15.5\%}
\providecommand{\acqQTwoIntAbs}{\$109.0 million}
\providecommand{\acqQTwoM}{34\%}
\providecommand{\acqQTwoMx}{34.1\%}
\providecommand{\acqQTwoOrg}{3\%}
\providecommand{\acqQTwoReorg}{\$51 million}
\providecommand{\acqQTwoReorgAbs}{\$50.8 million}
\providecommand{\acqQTwoReorgSh}{7.2\%}
\providecommand{\acqQTwoReorgX}{7.2\%}
\providecommand{\acqQTwoRev}{\$704.2 million}
\providecommand{\acqQTwoTrx}{0.9\%}
\providecommand{\acqRAdjQOne}{183.4}
\providecommand{\acqRAdjTwentyFive}{141.6}
\providecommand{\acqRAdjTwentyFour}{118.0}
\providecommand{\acqRAllInvQOne}{\ensuremath{-}129.3}
\providecommand{\acqRAllInvTwentyFive}{\ensuremath{-}24.8}
\providecommand{\acqRAllInvTwentyFour}{\ensuremath{-}28.7}
\providecommand{\acqRAllInvUQOne}{\ensuremath{-}117.5}
\providecommand{\acqRAllInvUTwentyFive}{\ensuremath{-}16.5}
\providecommand{\acqRAllInvUTwentyFour}{\ensuremath{-}25.0}
\providecommand{\acqRAnnHiTwentyFive}{92.8}
\providecommand{\acqRAnnHiTwentyFour}{86.8}
\providecommand{\acqRAnnLoTwentyFive}{84.4}
\providecommand{\acqRAnnLoTwentyFour}{80.9}
\providecommand{\acqRGaapQOne}{152.2}
\providecommand{\acqRGaapTwentyFive}{115.9}
\providecommand{\acqRGaapTwentyFour}{92.3}
\providecommand{\acqRIntQOne}{136.7}
\providecommand{\acqRIntTwentyFive}{105.0}
\providecommand{\acqRIntTwentyFour}{87.5}
\providecommand{\acqRLaAvgHiTwentyFive}{111.3}
\providecommand{\acqRLaAvgHiTwentyFour}{97.7}
\providecommand{\acqRLaAvgLoTwentyFive}{104.4}
\providecommand{\acqRLaAvgLoTwentyFour}{94.2}
\providecommand{\acqRLaFullHiTwentyFive}{101.1}
\providecommand{\acqRLaFullHiTwentyFour}{90.5}
\providecommand{\acqRLaFullLoTwentyFive}{92.7}
\providecommand{\acqRLaFullLoTwentyFour}{84.5}
\providecommand{\acqROcfQOne}{144.8}
\providecommand{\acqROcfSbcOnlyTwentyFive}{112.7}
\providecommand{\acqROcfSbcTwentyFive}{112.6}
\providecommand{\acqROcfSbcTwentyFour}{100.0}
\providecommand{\acqROcfTwentyFive}{116.9}
\providecommand{\acqROcfTwentyFour}{103.9}
\providecommand{\acqROcfUQOne}{156.6}
\providecommand{\acqROcfUTwentyFive}{125.2}
\providecommand{\acqROcfUTwentyFour}{107.6}
\providecommand{\acqROcfUZeroTwentyFive}{127.8}
\providecommand{\acqROrgIntNoReorgTwentyFive}{29.4}
\providecommand{\acqROrgIntNoReorgTwentyFour}{28.8}
\providecommand{\acqROrgIntTwentyFive}{23.4}
\providecommand{\acqROrgIntTwentyFiveR}{23}
\providecommand{\acqROrgIntTwentyFour}{21.1}
\providecommand{\acqROrgNoReorgTwentyFive}{40.3}
\providecommand{\acqROrgNoReorgTwentyFiveR}{40}
\providecommand{\acqROrgNoReorgTwentyFour}{33.7}
\providecommand{\acqROrgTwentyFive}{34.3}
\providecommand{\acqROrgTwentyFour}{26.0}
\providecommand{\acqRQTwoAdj}{180.5}
\providecommand{\acqRQTwoGap}{123}
\providecommand{\acqRQTwoOrg}{37}
\providecommand{\acqRQTwoOrgInt}{21.6}
\providecommand{\acqRQTwoOrgIntNoReorg}{28.9}
\providecommand{\acqRQTwoOrgIntNoReorgR}{29}
\providecommand{\acqRQTwoOrgIntR}{22}
\providecommand{\acqRQTwoOrgNoReorg}{44.3}
\providecommand{\acqRQTwoOrgNoReorgR}{44}
\providecommand{\acqRQTwoOrgX}{37.1}
\providecommand{\acqRQTwoRep}{160}
\providecommand{\acqRQTwoRepInt}{145.0}
\providecommand{\acqRQTwoRepIntR}{145}
\providecommand{\acqRQTwoRepX}{160.5}
\providecommand{\acqRUcAvgHiTwentyFive}{108.0}
\providecommand{\acqRUcAvgHiTwentyFour}{95.0}
\providecommand{\acqRUcAvgLoTwentyFive}{100.9}
\providecommand{\acqRUcAvgLoTwentyFour}{91.5}
\providecommand{\acqRUcEndHiTwentyFive}{95.2}
\providecommand{\acqRUcEndHiTwentyFour}{86.1}
\providecommand{\acqRUcEndLoTwentyFive}{87.2}
\providecommand{\acqRUcEndLoTwentyFour}{80.3}
\providecommand{\acqRdebt}{9.19\%}
\providecommand{\acqRdebtReg}{7.7\%}
\providecommand{\acqReminiFte}{30}
\providecommand{\acqReorgAbsTwentyFive}{\$78.6 million}
\providecommand{\acqReorgQOne}{12.6\%}
\providecommand{\acqReorgRatio}{3.6}
\providecommand{\acqReorgTwentyFive}{6.0\%}
\providecommand{\acqReorgTwentyFour}{7.7\%}
\providecommand{\acqReorgUsdEx}{\$30.1 million}
\providecommand{\acqReqCashHi}{246\%}
\providecommand{\acqReqCashLo}{124\%}
\providecommand{\acqReqHi}{76.0\%}
\providecommand{\acqReqLo}{71.7\%}
\providecommand{\acqReqOrg}{63.7\%}
\providecommand{\acqRetHi}{28\%}
\providecommand{\acqRetLo}{22\%}
\providecommand{\acqRetNetHi}{22\%}
\providecommand{\acqRetNetLo}{14\%}
\providecommand{\acqRetNetTauHi}{19.8\%}
\providecommand{\acqRetNetTauLo}{13.6\%}
\providecommand{\acqRetTauHi}{25.8\%}
\providecommand{\acqRetTauLo}{22.3\%}
\providecommand{\acqRevQOneFive}{\$258.9 million}
\providecommand{\acqRevQOneSix}{\$601.3 million}
\providecommand{\acqRevTwentyFive}{\$1{,}306.4 million}
\providecommand{\acqRevTwentyFour}{\$671.1 million}
\providecommand{\acqRevTwentyThree}{\$387.1 million}
\providecommand{\acqRpCustEight}{0.21\%}
\providecommand{\acqRpCustOne}{1.69\%}
\providecommand{\acqRpCustXbOne}{0.67\%}
\providecommand{\acqRpCustXbThree}{0.19\%}
\providecommand{\acqRpCustXbTwo}{0.34\%}
\providecommand{\acqRpFteDE}{387}
\providecommand{\acqRpFteFR}{705}
\providecommand{\acqRpFteUK}{1{,}124}
\providecommand{\acqRpFteUS}{337}
\providecommand{\acqRpGainBound}{11.40}
\providecommand{\acqRpGainDE}{6.45}
\providecommand{\acqRpGainFR}{5.87}
\providecommand{\acqRpGainHalf}{5.53}
\providecommand{\acqRpGainUS}{8.43}
\providecommand{\acqRpsQOne}{\$0.97 million}
\providecommand{\acqRpsTwentyFive}{\$2.57 million}
\providecommand{\acqRpsTwentyThree}{\$1.12 million}
\providecommand{\acqRtv}{8.93\%}
\providecommand{\acqRunRateInt}{572}
\providecommand{\acqRunRateIntHOne}{404}
\providecommand{\acqRunRateIntHi}{787}
\providecommand{\acqRunRateShPf}{14.7\%}
\providecommand{\acqRunRateShPfHi}{20.2\%}
\providecommand{\acqRunRateShQ}{20.3\%}
\providecommand{\acqRunRateShQHi}{27.9\%}
\providecommand{\acqRw}{8.93\%}
\providecommand{\acqRwAnnual}{9.34\%}
\providecommand{\acqSbcTwentyFive}{4.3\%}
\providecommand{\acqSbcTwentyFiveTwo}{4.25\%}
\providecommand{\acqSbcTwentyFour}{3.9\%}
\providecommand{\acqSbcTwentyThree}{3.8\%}
\providecommand{\acqSbcUsdTwentyFive}{\$55.5 million}
\providecommand{\acqScoopLoss}{\$84.5 million}
\providecommand{\acqSpo}{621}
\providecommand{\acqSpoAvgTwentyFive}{508}
\providecommand{\acqSpoAvgTwentyThree}{346}
\providecommand{\acqSpoShare}{27\%}
\providecommand{\acqSpoTech}{42\%}
\providecommand{\acqStandCashHi}{30.5\%}
\providecommand{\acqStandCashLo}{8.9\%}
\providecommand{\acqStandHi}{5.1\%}
\providecommand{\acqStandLo}{9.4\%}
\providecommand{\acqStandUsdHi}{\$67 million}
\providecommand{\acqStandUsdLo}{\$123 million}
\providecommand{\acqSubsTwentyFive}{93\%}
\providecommand{\acqSubsTwentyFour}{92\%}
\providecommand{\acqSubsTwentyThree}{95\%}
\providecommand{\acqTalent}{55}
\providecommand{\acqTalentShare}{9\%}
\providecommand{\acqTargets}{200}
\providecommand{\acqTauIres}{24\%}
\providecommand{\acqTenFive}{48\%}
\providecommand{\acqTenMean}{8.0}
\providecommand{\acqTenSqFive}{23\%}
\providecommand{\acqTenTen}{28\%}
\providecommand{\acqTirr}{5}
\providecommand{\acqTlbEur}{EUR 395 million}
\providecommand{\acqTlbUsd}{\$1.25 billion}
\providecommand{\acqTotalAssets}{\$6.98 billion}
\providecommand{\acqTpBe}{13.5\%}
\providecommand{\acqTpBeAnn}{14.4\%}
\providecommand{\acqTpIl}{10.2\%}
\providecommand{\acqTpIm}{10.1\%}
\providecommand{\acqTpImAnn}{10.6\%}
\providecommand{\acqTpIp}{9.1\%}
\providecommand{\acqTpIuSp}{12.4\%}
\providecommand{\acqTpMult}{18.2}
\providecommand{\acqTpPmaxM}{10.6}
\providecommand{\acqTpPmaxP}{4.1}
\providecommand{\acqTpPmaxShare}{54\%}
\providecommand{\acqTpPrice}{18.5}
\providecommand{\acqTpUplift}{106\%}
\providecommand{\acqTpUpliftFast}{101\%}
\providecommand{\acqTractive}{\$759 million}
\providecommand{\acqTractiveWiki}{\$900 million}
\providecommand{\acqTrep}{5}
\providecommand{\acqTrxQOne}{6.3\%}
\providecommand{\acqTrxTwentyFive}{6.5\%}
\providecommand{\acqTrxTwentyFour}{4.8\%}
\providecommand{\acqTtmNi}{\$139.5 million}
\providecommand{\acqTtmRev}{\$1.65 billion}
\providecommand{\acqUcPtsAnnHi}{0.30}
\providecommand{\acqUcPtsHi}{0.59}
\providecommand{\acqUnlForLMax}{31\%}
\providecommand{\acqUnlForLMin}{30\%}
\providecommand{\acqUnpricedDebt}{2{,}234}
\providecommand{\acqVimDays}{327}
\providecommand{\acqVimPre}{\$376.9 million}
\providecommand{\acqVimeoB}{\$1.38 billion}
\providecommand{\acqVimeoPrem}{91\%}
\providecommand{\acqVimeoPx}{\$7.85}
\providecommand{\acqWetPlanned}{75\%}
\providecommand{\acqWetStaff}{350}
\providecommand{\acqXbOne}{2.5}
\providecommand{\acqXbThree}{9}
\providecommand{\acqXbTwo}{5}
\providecommand{\acqXgRatio}{0.51}
% permanent-capital row of tab:reg-presets (generated by code/verify_part2_acquirer.py; \providecommand, so a second input is harmless)

Sections~\ref{sec:tasks} and \ref{sec:ext} value transitions at one discount rate over a going-concern horizon. The owner changes that: the rate, the date at which it stops counting, the rule that turns savings into a sale price, the pressure on reported earnings, the financing constraints on the transition dip and the speed of decisions. This section makes the ownership regime an input. The terminal rule matters as much as the rate, and rankings of activities flip between regimes (Proposition~\ref{prop:reg-value}). Each regime acts on its own Rule-of-$X$ weight, and a sponsor's is close to investors' screening weights (Proposition~\ref{prop:reg-weight}). A hurdle above the cost of capital raises every breakpoint, while a short hold caps the time to run-rate rather than the payback unless buyers refuse to pay for recent savings (Proposition~\ref{prop:reg-speed}). Covenants and earnings pressure turn the dip of Proposition~\ref{prop:tasks-scorepath} into a budget under which the program ramps up instead of front-loading (Proposition~\ref{prop:reg-dip}). A take-private pays for a transformation only when public ownership would block most of it (Proposition~\ref{prop:reg-lbo}). Stakeholder preferences are parameters, on which we take no position, and every example is illustrative, conditional on presets marked as assumed in Table~\ref{tab:reg-presets}. Proofs are in Appendix~\ref{app:proofs-regimes}; \texttt{code/verify\_part2\_regimes.py} checks every closed form and number.

\subsection{Regimes}\label{sec:regimes-setup}

\begin{assumption}[Ownership regime]\label{ass:reg-regime}
A regime $\mathsf o\in\{\mathrm{pub},\mathrm{fam},\mathrm{pe},\mathrm{vc}\}$ (public, founder or family, private-equity-backed, venture-backed) fixes six things.
\begin{enumerate}
\item[(i)] (Rate and horizon.) A rate $r_{\mathsf o}>0$ and a horizon $T_{\mathsf o}\in(0,\infty]$: the cost of capital $r_w$ and $T=\infty$ (public); $r_{\rm fam}$ and a long horizon (family); the sponsor's target gross return $r_{\rm pe}$ and remaining hold; $r_{\rm vc}$ and $T_{\rm vc}$. Rates are continuously compounded; an annual preset enters as $r=\ln(1+r_{\rm annual})$.
\item[(ii)] (Terminal rule.) At a finite horizon the owner sells. A change in run-rate EBITDA $\Delta\mathrm{EB}$ or free cash flow $\Delta W$ in place for the seasoning period $T_{\rm sea}$ adds $h_xM_{\rm exit}\Delta\mathrm{EB}$ to the price under the \emph{EBITDA rule} (the buyer pays $M_{\rm exit}$ times EBITDA on the existing business and capitalizes the share $h_x\in[0,1]$ of new seasoned run-rate), and $\mathfrak A_T\Delta W$ under the \emph{market rule}, where $\mathfrak A_T:=\mathcal A_T(\eta)+\theta\eta^T\Gamma$ is the value of a permanent dollar of free cash flow in Lemma~\ref{lem:valuation2}; there a change $\Delta g$ in exit growth adds $\mathfrak A_TX^{\rm scr}S\,\Delta g$. Unseasoned run-rate adds nothing.
\item[(iii)] (Short-term pressure.) A weight $\mathfrak w_{\mathsf o}\ge0$ on shortfalls in reported operating income: the owner's objective falls by $\mathfrak w_{\mathsf o}\sum_{y=1}^{Y_m}e^{-r_{\mathsf o}y}(-\Delta\mathrm{OI}_y-\mathfrak gS)^+$, with $\Delta\mathrm{OI}_y$ the change in GAAP operating income in year $y$, $Y_m$ the window and $\mathfrak g\ge0$ the dip guided in advance, as a share of revenue $S$.
\item[(iv)] (Financing.) A sponsor's lenders require $\mathrm{ND}_t\le\overline{\rm lev}\,\mathrm{EB}_t$, with $\mathrm{ND}_0={\rm lev}_0\mathrm{EB}_0$, ${\rm lev}_0<\overline{\rm lev}$; the current wave's flow enters tested EBITDA with weight $w_{\rm cov}\in(0,1]$ (one if annualized, about $T_{\rm par}$ on a trailing twelve months). Other firms face a cash budget, net of required dividends and buybacks.
\item[(v)] (Stakeholder terms.) Employment preference is the altruism parameter of Assumption~\ref{ass:ext-theta}(iv), $\omega_h=(1-\tau)b_{\rm fam}w_j$, so the owner counts a displaced wage dollar at $1-b_{\rm fam}$ and a transition's saving at $W-b_{\rm fam}(1-u_d)C^{\rm disp}$. Customers receive the pass-through share $\bar\varphi$ (Lemma~\ref{lem:ext-retained}); equity pay enters $w_j$ at fair value.
\item[(vi)] (Decision speed.) The decision latency $\ell_{\rm dec}$ of Assumption~\ref{ass:ext-bandwidth} becomes ${\rm lat}_{\mathsf o}\ell_{\rm dec}$.
\end{enumerate}
\end{assumption}

\begin{definition}[Regime value]\label{def:reg-value}
For a plan with expected incremental cash flows $\mathrm{CF}$, terminal rule $\mathrm{TV}_{\mathsf o}$ of Assumption~\ref{ass:reg-regime}(ii) and planned GAAP path $\Delta\mathrm{OI}_y$, the regime value is
\begin{equation}\label{eq:reg-value}
V_{\mathsf o}:=\int_{[0,T_{\mathsf o})}e^{-r_{\mathsf o}t}\,d\,\mathbb E\,\mathrm{CF}(t)+e^{-r_{\mathsf o}T_{\mathsf o}}\,\mathbb E\,\mathrm{TV}_{\mathsf o}-\mathfrak w_{\mathsf o}\sum_{y=1}^{Y_m}e^{-r_{\mathsf o}y}\bigl(-\Delta\mathrm{OI}_y-\mathfrak gS\bigr)^+,
\end{equation}
with no terminal term if $T_{\mathsf o}=\infty$.
\end{definition}

Table~\ref{tab:reg-presets} lists the presets. Private equity's returns, holds and fees are documented \citep{kaplan2005private,harris2014private,kaplan2009leveraged}, and the median buyout borrows 5.2 times EBITDA \citep{axelson2013borrow}. Hurdles exist only as stated targets: an average of 22\% for buyout investors \citep{gompers2016what} and 31\% for venture investors \citep{gompers2020how}, both surveys. The cost of capital is a practitioner data set \citep{damodaran2026cost} and the exit multiple a small practitioner sample \citep{silverpeak2026software}. The myopia weight has a direction \citep{stein1989efficient,graham2005economic,asker2015corporate} but no estimate; covenant levels have a supervisory precedent and documented consequences of violation \citep{chava2008how}; family-firm parameters are assumed \citep{villalonga2006how}. The permanent-capital row is a serial acquirer's disclosed policy \citep{bendingspoons2026f1} and the coupon on its first secured notes \citep{bendingspoons2026sixk}, analyzed in Section~\ref{sec:acq-pc}; the sponsor's debt rate stays at its preset. The window $Y_m$ and the guided slack $\mathfrak g$ are scenario parameters: the examples use $Y_m=2$ years, $\mathfrak g=0$ for a single transition and $\mathfrak g=0.01$ (one point of revenue) for a program.

\emph{Case R.} Several examples use one transition of Assumption~\ref{ass:tasks-transition}, independent of the running example of Section~\ref{sec:tasks} (all illustrative, in \$ thousands): a process with legacy cost 1{,}000 and in-house labor-equivalent spend 1{,}200 a year moves to an AI mode that saves $W=400$ a year for $T_h=10$ years, with $F=1{,}000$, $I=0$ and $C^J=50$; its parallel run lasts $T_{\rm par}=0.17$ years (the parameter file's preset for an AI agent in shadow mode) with $\omega_{\rm par}=0.8$ and $\nu_{\rm rec}=0.1$, so $C^{\rm par}=920$ a year; $C^{\rm str}=150$, $C^{\rm wo}=50$, $L_{\rm roll}=500$, $\varpi_p=0.8$, $\bar b_p=0.67$, $T_{\rm det}=0.25$ and $T_A=4$. $T_{\rm par}$ and $T_A$ are grounded presets, and $\bar b_p$ is the engine's realized share, the share of targeted value that self-reported successful transformations realize; the rest are assumed.

\begin{table}[t]
\centering
\caption{Regime presets. Status: S peer-reviewed source, Sp practitioner data, Ss survey, Sc secondary statement, F company filing, A assumed. Sources and quality flags are recorded in \texttt{data/tasks/params.yaml}, block \texttt{ownership\_regimes}}\label{tab:reg-presets}
\scriptsize
\setlength{\tabcolsep}{3pt}
\begin{tabular}{@{}p{2.9cm}p{3.1cm}p{2.6cm}p{4.0cm}p{2.5cm}@{}}
\toprule
Parameter & Public & Family & PE-backed & VC-backed\\
\midrule
Rate $r_{\mathsf o}$, annual (continuous) & $r_w$: 9.34\% (8.93\%), range 7.0 to 10.7, Sp & 8.34\% (8.01\%), $r_w-1$ pt, A & 22\% (19.9\%), 20 to 25, Ss & 31\% (27.0\%), 28 to 35, Ss\\
Horizon $T_{\mathsf o}$ & $\infty$, A & $\infty$ (file: 30 yr, 10 to 100), A & hold 6 yr (4 to 7), S \citep{kaplan2009leveraged} & 7 yr (5 to 10), A\\
Terminal rule & none & none & EBITDA, $M_{\rm exit}=14$ (10 to 17; band 12 to 17), Sp; $h_x=1$ (0 to 1), $T_{\rm sea}=1$ yr, A & market, $\mathfrak A_T$ and $X^{\rm scr}$ (Part I)\\
Myopia $\mathfrak w$; window $Y_m$, slack $\mathfrak g$ & 0.30 (0 to 0.6), A; dual class 0.15, A; $Y_m$, $\mathfrak g$ scenario & 0 & 0 & 0\\
Leverage ${\rm lev}_0$, covenant $\overline{\rm lev}$ & & & 5.2 total debt (4.2 to 6.1), S; 7.0 (6 to 8), Sc floor, central A; $w_{\rm cov}=1$, debt 8\% (7.7\%), 6 to 11, A & \\
Premium $\mathfrak p^{\rm req}$; fees, carry, preferred return & 30\% (15 to 50), S range, central A & & 2\%, 20\%, S \citep{kaplan2009leveraged}; 8\% (6 to 8), Sc & \\
Employment preference $b_{\rm fam}$ & 0 & 0.10 (0 to 0.3), A & 0 & 0\\
Latency factor ${\rm lat}_{\mathsf o}$ & 1.0, A & 0.7, A & 0.5, A & not set\\
\midrule
Permanent capital, serial acquirer (Assumption~\ref{ass:acq-pc}, Proposition~\ref{prop:acq-pc}) & \multicolumn{4}{p{12.8cm}@{}}{Horizon $\infty$, no exit, F. Rate: hurdles on the method IRR \eqref{eq:acq-method}, $\bar{\mathfrak i}_{\rm U}=25\%$ (22.3\%) unlevered and $\bar{\mathfrak i}_{\rm L}=65\%$ (50.1\%) levered, F, with flows after $T^{\rm IRR}=5$ yr valued at $r_{\rm tv}=r_w$, A. Debt $\mathrm{ND}_0=\min\{\overline{\mathrm{LTV}}P,\mathrm{ND}^{\rm rep}\}$ with $\overline{\mathrm{LTV}}=85\%$ of the enterprise value paid and $\mathrm{ND}^{\rm rep}$ the debt that five years of free cash flow repays, F; covenant $\mathrm{ND}_t\le4.0\,\mathrm{EB}_t$, F; debt rate $\ln(1+\acqNotesUsdC)=\acqRdebt$, the coupon on the \acqNotesUsd{} senior secured notes priced on 28 October 2026 (\acqNotesEurC{} on the \acqNotesEur{} tranche), F \citep{bendingspoons2026sixk}. Myopia $\mathfrak w_{\rm pc}\in[0,\mathfrak w_{\rm dc}]$ (private, then a dual-class listing), A. Capability $\mathcal U_{\rm pc}$ of Assumption~\ref{ass:acq-owner}. Filing: \citet{bendingspoons2026f1}.}\\
\bottomrule
\end{tabular}

\smallskip
\begin{minipage}{0.97\linewidth}\scriptsize
Notes: annual presets, continuous equivalents $\ln(1+r)$ in parentheses (used in every example; IRRs are reported annually compounded, and the preferred return enters as $1.08^T$). The covenant floor is a supervisory concern level of 6 times; the leverage median is total debt (net debt 5.1). The exit band is a practitioner report with restricted republication, accessed 29 October 2026. Examples set the entry multiple equal to the exit multiple, so no IRR includes multiple expansion (the file's historical all-industry entry multiple is not paired with the 2025 software exit multiple).
\end{minipage}
\end{table}

\subsection{The value of a transition by regime}\label{sec:regimes-value}

\begin{proposition}[Regime value of a transition]\label{prop:reg-value}
Consider the transition of Assumption~\ref{ass:tasks-transition}, let $W^{\rm E}:=W+\delta_KI$ be its run-rate EBITDA saving (the cash saving plus the replacement capital spending it requires, which EBITDA does not charge), and suppose cutover and seasoning end before the horizon, $T_{\rm par}+T_{\rm sea}\le T_{\mathsf o}$.
\begin{enumerate}
\item[(a)] (Value.) The regime value is
\begin{equation}\label{eq:reg-trans}
V_{\mathsf o}=\mathrm{NPV}_p\big|_{r=r_{\mathsf o},\,T_h=\min\{T_h,T_{\mathsf o}\}}+e^{-r_{\mathsf o}T_{\mathsf o}}\varpi_p\bar b_p\,\widetilde M\,W-\mathcal P_{\mathsf o},
\end{equation}
where $\mathrm{NPV}_p$ is \eqref{eq:tasks-npv}, $\widetilde M=h_xM_{\rm exit}W^{\rm E}/W$ under the EBITDA rule, $\widetilde M=\mathfrak A_T$ under the market rule, $\widetilde M=0$ for a going concern, and $\mathcal P_{\mathsf o}$ is the myopia term of \eqref{eq:reg-value} on the planned GAAP path of Proposition~\ref{prop:tasks-scorepath}(b).
\item[(b)] (Capitalization of a run-rate dollar.) For a step transition (cost $C$ at date 0, saving $W$ a year from $t$ on), $V_{\mathsf o}=-C+W\mathfrak K_{\mathsf o}(t)$, with $\mathfrak K^{\rm gc}(t)=e^{-rt}/r$ for a going concern and
\begin{equation}\label{eq:reg-K}
\mathfrak K^{\rm x}(t)=\frac{e^{-rt}-e^{-rT}}{r}+e^{-rT}\widetilde M\ \ (t\le T-T_{\rm sea}),\qquad \mathfrak K^{\rm x}(t)=\frac{(e^{-rt}-e^{-rT})^+}{r}\ \ (t>T-T_{\rm sea})
\end{equation}
under an exit rule. For $0<t\le T-T_{\rm sea}$, $\mathfrak K^{\rm x}(t)/\mathfrak K^{\rm x}(0)>e^{-rt}$ if and only if $\widetilde M>1/r$: within the hold, an exit multiple above $1/r$ makes the owner more patient than a going concern at the same rate, and the patience ends at $T-T_{\rm sea}$.
\item[(c)] (Ranking flips.) Let a quick design $q$ and a slow design $\ell$ of the same process have $C_\ell\ge C_q$, $t_\ell>t_q$ and $W_\ell>W_q$. For a going concern, $re^{rt_q}(V_\ell-V_q)$ is strictly decreasing in $r$, so there is a unique rate $r^\ddagger$ above which the quick design is preferred. Under an exit rule, the slow design loses its terminal value if $t_\ell>T-T_{\rm sea}$.
\item[(d)] (EBITDA versus cash.) A capital-for-labor switch of Proposition~\ref{prop:tasks-capital} (capital $\Delta K_j$ per unit of annual volume, operating saving $\Delta c_j$ per unit, constant volume) raises the value of an owner with the EBITDA rule if and only if
\begin{equation}\label{eq:reg-zeta}
\frac{\Delta c_j}{\Delta K_j}>\zeta^{\rm pe}:=\frac{1+\delta_K\mathcal A_r(T)}{\mathcal A_r(T)+e^{-rT}h_xM_{\rm exit}},\qquad \mathcal A_r(T):=\frac{1-e^{-rT}}{r},
\end{equation}
at $r=r_{\rm pe}$, $T=T_{\rm pe}$, against $\Delta c_j/\Delta K_j>r_w+\delta_K$ for a public going concern. If $\zeta^{\rm pe}<r_w+\delta_K$, switches between the two hurdles are adopted by the sponsor and rejected on value.
\end{enumerate}
\end{proposition}

\noindent\emph{Intuition.} An owner who sells is paid for run-rate at exit, not for savings it will not collect. If the buyer capitalizes seasoned run-rate at more than $1/r_{\mathsf o}$, when a saving starts matters less than to a going concern at the same rate, provided it seasons before exit; after that date it matters completely. Patience depends on the terminal rule, not only the rate. Part (d) follows from capitalizing EBITDA, which treats replacement capital as free, so the sponsor's hurdle on capital can be lower than the public owner's although its discount rate is higher. Buyout gains in operations and returns that track industry multiples \citep{guo2011do} fit both channels, and markets capitalized workforce exposure to generative AI within weeks of its release \citep[a working paper]{eisfeldt2023generative}, as the market rule does.

\begin{table}[t]
\centering
\caption{Regime values of three transitions (\$ millions, illustrative values, conditional on assumed presets)}\label{tab:reg-rank}
\footnotesize
\begin{tabular}{@{}lccccccccc@{}}
\toprule
Transition & $C$ & $W$ ($W^{\rm E}$) & $t$ & Public & Family & PE, $h_x=1$ & PE, $h_x=0.5$ & VC\\
\midrule
Q: quick cost-out & 2.0 & 1.0 (1.0) & 0.5 & 8.71 & 8.55 & 5.27 & 3.15 & 7.93\\
P: data platform, end-to-end automation & 6.0 & 2.0 (2.0) & 4.0 & 9.67 & 9.95 & 3.98 & $-0.26$ & 9.90\\
K: owned inference capacity & 4.0 & 0.3 (1.3) & 0.5 & $-0.79$ & $-0.40$ & 2.43 & $-0.33$ & $-1.02$\\
\bottomrule
\end{tabular}

\smallskip
\begin{minipage}{0.95\linewidth}\footnotesize
Notes: continuous rates. Public $r_w=0.0893$, $T=\infty$; family $r=0.0801$, $T=\infty$, $b_{\rm fam}=0.10$, displaced wage bills 1.2, 2.4 and 0 a year (the parameter file's 30-year horizon, without terminal value, gives Q 7.56 and P 7.96 and keeps the ranking); PE $r=0.199$, a fresh six-year hold, $M_{\rm exit}=14$, $h_x\in\{1,0.5\}$; VC $r=0.270$, $T=7$, market rule with $\mathfrak A_T=48.0$, the value of Lemma~\ref{lem:valuation2} at Part I's discussion parameters ($r=0.10$, $g=0.25$, $T=10$, $G_\infty=0.03$, $\theta=1/2$). Savings constant, run-rate from $t$. Values computed by \texttt{code/verify\_part2\_regimes.py}.
\end{minipage}
\end{table}

\noindent\emph{Example (illustrative).} In Table~\ref{tab:reg-rank} the public, family and venture owners rank the slow platform P above the quick cost-out Q, and the sponsor ranks Q first. As alternative designs, Q and P switch at $r^\ddagger=9.9\%$, about one point above the public cost of capital; the family's assumed lower rate widens P's lead despite its larger displaced wage bill. The venture owner discounts at 27\% (31\% annual) yet ranks like the public firm, because the market rule capitalizes savings at a growth firm's margin multiple. Owned inference capacity K loses value for every owner except a sponsor whose buyer capitalizes all seasoned run-rate. With $\delta_K=1/3$ and $h_x=1$, \eqref{eq:reg-zeta} gives $\zeta^{\rm pe}=0.28$ per dollar of capital against $r_w+\delta_K=0.42$, so a capital-for-labor switch that saves \$0.60 per unit on \$1.50 of capital per unit of annual capacity is worth $+\$168{,}000$ to the sponsor on 120{,}000 units a year and $-\$46{,}000$ at the public cost of capital. These flips depend on the assumed $h_x$: at $h_x=0.5$, $\zeta^{\rm pe}=0.385$, so that switch (0.40) still pays the sponsor (\$15{,}000), but K is worth $-0.33$ and P $-0.26$. All values in this section are before tax. With full expensing of domestic server capital ($\psi=1$; Section~\ref{sec:acct-standards}), the going-concern capital tests $\Delta c_j>(r_{\mathsf o}+\delta_K)\Delta K_j$ hold after tax unchanged, since savings and capital charges both scale by $1-\tau$; for servers used predominantly abroad, or after an election out, the capital charge is multiplied by $\Upsilon=(1-\psi\tau)/(1-\tau)>1$ (Section~\ref{sec:ext-setup}), and the sponsor's $\zeta^{\rm pe}$, which capitalizes EBITDA before tax, is a pre-tax hurdle. For a public owner, Case R cuts first-year GAAP operating income by \$290{,}000 (Proposition~\ref{prop:tasks-scorepath}(c)); with no guided slack and $\mathfrak w=0.3$ this costs \$80{,}000.

\subsection{Mapping to Part I: the regime's growth weight}\label{sec:regimes-partI}

\begin{proposition}[The regime's growth weight]\label{prop:reg-weight}
Let $\bar\sigma=\infty$.
\begin{enumerate}
\item[(a)] (Going concern.) An owner with rate $r_{\mathsf o}>G$, an infinite horizon and $\mathfrak w_{\mathsf o}=0$ maximizes $R_X$ at every date with $X_{\mathsf o}=v'(s;r_{\mathsf o})$, Proposition~\ref{prop:dynamic} at $r=r_{\mathsf o}$. In the small-share limit $X_{\mathsf o}\to v_0(r_{\mathsf o})$, which is strictly decreasing in $r_{\mathsf o}$.
\item[(b)] (Short-term pressure.) If each dollar of current GAAP operating income $m_{op}S$ carries the extra weight $\mathfrak w$ (growth spend is expensed, growth capital is not), then at a given weight $X$ the per-date maximand is $(1+\mathfrak w)R_{X^{\rm my}}$ with
\begin{equation}\label{eq:reg-Xmy}
X^{\rm my}:=k+\frac{X-k}{1+\mathfrak w}.
\end{equation}
Excess growth falls by the factor $\rho:=(1+\mathfrak w)^{-\alpha/(1-\alpha)}$ and growth spend by $(1+\mathfrak w)^{-1/(1-\alpha)}$, and the loss in the true criterion $R_X$ is $\Pi(X)\bigl[1-(\rho-\alpha\rho^{1/\alpha})/(1-\alpha)\bigr]$.
\item[(c)] (Hold and exit.) For any horizon and terminal value, the per-date maximand is $R_X$ with $X=\partial V/\partial S$. If a marginal dollar of revenue is permanent, grows at $G$, earns the free-cash-flow margin $\hat m$ and is worth $X_T$ at exit, then with $\tau$ years to exit
\begin{equation}\label{eq:reg-Xexit}
X^{\rm x}(\tau)=\bigl(1-e^{-(r_{\mathsf o}-G)\tau}\bigr)X^{\rm G}+e^{-(r_{\mathsf o}-G)\tau}X_T,\qquad X^{\rm G}:=\frac{\hat m}{r_{\mathsf o}-G},
\end{equation}
with $X_T=h_xM_{\rm exit}\hat m^{\rm E}$ under the EBITDA rule ($\hat m^{\rm E}$ the marginal EBITDA margin) and $X_T=(V/S)_T$ under the market rule. $X^{\rm x}$ moves monotonically toward $X_T$ as exit approaches.
\item[(d)] (Frontier.) Because $\Delta^*$ is increasing in $X$ (Proposition~\ref{prop:score}(d)), regimes with higher weights choose faster growth, more spend and lower margins, and $R^*_{X_{\mathsf o}}=m_0+(X_{\mathsf o}-k)G+C_\alpha[\ell(X_{\mathsf o}-k)]^{1/(1-\alpha)}$.
\item[(e)] (Growth engine or cost-out.) A proportional saving $\delta_\sigma$ on growth activities raises $R^*_X$ by $\Pi(X)[(1-\delta_\sigma)^{-\alpha/(1-\alpha)}-1]$ (Proposition~\ref{prop:tasks-aggregate}(b)) and a cut $da$ in non-discretionary cost raises it by $da$. The growth engine ranks first if and only if
\begin{equation}\label{eq:reg-Xdag}
X>X^\ddagger:=k+\frac1\ell\Bigl[\frac{da}{C_\alpha\bigl((1-\delta_\sigma)^{-\alpha/(1-\alpha)}-1\bigr)}\Bigr]^{1-\alpha}.
\end{equation}
\end{enumerate}
\end{proposition}

\noindent\emph{Intuition.} Every owner solves Part I's per-date problem with its own price of growth. A lower rate raises the marginal value of revenue; a weight on reported earnings divides the net weight $X-k$ by $1+\mathfrak w$, Stein's underinvestment \citep{stein1989efficient} as a lower Rule-of-$X$ weight; managing to forecasts goes with lower R\&D growth and, in an estimated model, slows growth \citep{terry2023macro}. A sponsor values revenue for its cash until exit and the multiple at exit, both at its hurdle, so its weight lands near the 2 to 3 of investors' screens: both capitalize a few years of margin, not a permanent stream. Surveyed buyout investors say they add value more through growth than through cost cuts \citep{gompers2016what}; in the model that needs an exit price that rewards growth. Part (e) makes Part I's priorities (Corollary~\ref{cor:priorities}(ii)) a regime test; firm AI investment has gone with faster growth through product innovation \citep{babina2024artificial}, the channel a high weight rewards.

\noindent\emph{Example (illustrative).} Take $m_0=0.35$, $G=0.03$, $k=0$, $\alpha=0.3$ and demand pull $\ell=0.0535$, which makes the model's going-concern weight $v_0=8$ at $r_w=8.93\%$; for this firm it is at least $m_0/(r_w-G)=\regVzeroMin$ at any demand pull. At the assumed family rate of 8.01\% it is 10.8, so the weight is sensitive to the rate; the assumed $\mathfrak w=0.3$ lowers the public weight to 6.15. A sponsor with $\hat m=\hat m^{\rm E}=0.25$, $r_{\rm pe}=19.9\%$, $M_{\rm exit}=14$ and $h_x=1$ has $X^{\rm x}=2.21$, 2.70 and 3.50 at six, three and zero years to exit. At the mid-hold weight, the four owners grow at 5.5, 5.2, 5.0 and 4.4\% (family, public, pressured public, sponsor) with FCF margins of 26.8, 29.7, 31.3 and 33.9\%. A 20\% saving on growth activities beats a one-point overhead cut if and only if $X>6.85$: at the model's weights it is worth 1.25 points of score to the public firm, 0.86 to the pressured public firm and 0.26 to the sponsor, so the assumed earnings pressure reverses the public firm's ranking. Myopia costs the pressured firm 0.22 points of its true criterion a year. These weights are model outputs at assumed parameters; the data do not identify a firm's weight (Part~I), so the ranking is better read across owner weights at the same demand pull. At illustrative weights of \regXbOne, \regXbTwo{} and \regXbThree{} the growth-engine saving is worth \regXbGainOne, \regXbGainTwo{} and \regXbGainThree{} points against the one point of the overhead cut, so the overhead cut ranks first at the first two for every owner, and at \regXbThree{} the growth engine ranks first for the public owner and also for the pressured one, whose weight is \regXbMyThree. Earnings pressure reverses the public ranking only for weights between \regXbBandLo{} and \regXbBandHi, which the model's 8 is in and none of the three is. The sponsor ranks the overhead cut first at every weight; its weight is below the public owner's at \regXbTwo{} and \regXbThree, and at \regXbOne{} it is above it from three years before exit. % XBOUNDS: placeholder weights (macros regXb*, set in code/verify_part2_regimes.py)

\subsection{Speed: the pace of transformation by regime}\label{sec:regimes-speed}

\begin{proposition}[Pace by regime]\label{prop:reg-speed}
\begin{enumerate}
\item[(a)] (A hurdle raises every breakpoint.) The capital recovery factor $a(r,T)=r/(1-e^{-rT})$ is strictly increasing in $r$. If the alternatives carry no capital, a higher $r_{\mathsf o}$ raises $v^*_j$ and the make-or-buy volume, lowers $\pi^*_j$, tightens the capital test to $\Delta c_j>(r_{\mathsf o}+\delta_K)\Delta K_j$ and does not advance the adoption date of Proposition~\ref{prop:ext-timing}(b). A transition whose expected net cash flow changes sign once, from negative to positive (for Assumption~\ref{ass:tasks-transition}, when $C^{\rm par}\ge\varpi_p\bar b_pW$), creates value for a going concern if and only if its internal rate of return exceeds $r_{\mathsf o}$, and its minimum horizon $T^*_h(r_{\mathsf o})$ is nondecreasing.
\item[(b)] (What a hold caps.) Under an exit rule, let a step transition with time to run-rate $t_{\rm run}$ start at $t_0$, with $\tau=T-t_0$ left. It earns terminal value only if $t_0+t_{\rm run}+T_{\rm sea}\le T$. Its value is monotone in $t_0$, with derivative $e^{-rt_0}(rC-We^{-rt_{\rm run}})$, so the optimal start is immediate if the cash yield $We^{-rt_{\rm run}}/C$ is at least $r_{\mathsf o}$, and otherwise the latest admissible date $T-t_{\rm run}-T_{\rm sea}$ (or never, if the value there is negative). If buyers capitalize nothing ($h_x=0$), a fundable transition has simple payback $C/W\le\tau-t_{\rm run}$; with $h_x\widetilde M>0$ the admissible payback is $[e^{-rt_{\rm run}}-e^{-r\tau}]/r+e^{-r\tau}\widetilde M$, which can far exceed the remaining hold.
\item[(c)] (Program pace.) In the wave model of Proposition~\ref{prop:ext-waves}, under an exit rule the terminal value of a seasoned process does not depend on when it was migrated, so within the admissible window delay costs only interim savings, and the program solves \eqref{eq:ext-progcost} with the deadline $n_w\le\bar n_w:=\lfloor(T_{\mathsf o}-T_{\rm sea})/T_{\rm par}\rfloor$. Decision speed enters through the load breakpoint of Proposition~\ref{prop:ext-bandwidth}(a), $\hat\varkappa=1+c/2-\sqrt{c+c^2/4}$ with $c=r_d\ell_0{\rm lat}_{\mathsf o}$, which is strictly decreasing in $c$, so lower latency raises the bound on the pace in Proposition~\ref{prop:ext-bandwidth}(b).
\end{enumerate}
\end{proposition}

\noindent\emph{Intuition.} A higher rate charges builds and capital more heavily, so every threshold of Section~\ref{sec:tasks-breakpoints} with a fixed or capital cost moves against automation. A hold acts differently: the sponsor wants each transition seasoned before exit, and within that window its only reason to hurry is interim cash, which it compares with its hurdle. A short hold caps the payback each wave can have only when buyers do not pay for recent run-rate; when they do, it caps the time to run-rate plus seasoning. A sponsor's faster cost-out is a model implication, not documented behavior: in a large hospital buyout, margins rose through staffing and vendors with little new technology \citep{richards2024managing}, and we know of no verified study of AI adoption by ownership regime.

\noindent\emph{Example (illustrative).} With $T_A=4$ and no capital in mode A, the volume threshold of Proposition~\ref{prop:tasks-modes}(e) is 18.9\% higher at $r_w$ and 45.0\% higher at the sponsor's hurdle than with straight-line amortization. Case R has an internal rate of return of 9.1\% (9.5\% annually), barely above the public cost of capital: it is worth $+\$12{,}000$ at $r_w$ and $-\$471{,}000$ at the hurdle, and its minimum horizon is 9.9 years at $r_w$, just inside its ten-year life, and infinite at the hurdle; a permanent saving would earn 15.3\%. For the sponsor ($h_x=1$, $T_{\rm sea}=1$ year, both assumed), Q has cash yield 0.45, above the 0.199 hurdle, and should start now, while P yields 0.15 and is worth 4.24 if started a year later against 3.98 now. If buyers capitalized nothing, Q would still be worth 1.03 and P $-4.51$. With 1.5 years left and one year of seasoning, the twelve-process program of Proposition~\ref{prop:ext-waves} must run as one wave, at a cost of 25.7 against 19.7 for its three optimal waves. Latency factors of 1.0, 0.7 and 0.5 give load breakpoints of 0.70, 0.74 and 0.78.

\subsection{The dip budget and the constrained wave structure}\label{sec:regimes-dip}

\begin{assumption}[Dip budgets in the wave model]\label{ass:reg-dip}
The program of Assumption~\ref{ass:ext-program} has identical processes, each with a one-time cash cost $\mathcal T_1$ at the start of its wave. During wave $k$ the margin flow changes by $sN_{k-1}-D(n_k)$ a year, where $N_{k-1}:=\sum_{i<k}n_i$ (Proposition~\ref{prop:ext-waves}(f)). A public firm guides a dip of $\mathfrak gS$ a year and holds to it (hard budget) or pays $\mathfrak w T_{\rm par}(D(n_k)-sN_{k-1}-\mathfrak gS)^+$ for wave $k$ (soft budget). A sponsor's covenant counts one-time costs in net debt and the current wave's flow in EBITDA, omitting the cash of past waves' dual running (exact if it is financed outside the tested debt, first order in $T_{\rm par}/(w_{\rm cov}\overline{\rm lev})$ otherwise).
\end{assumption}

\begin{proposition}[Dip budget and constrained waves]\label{prop:reg-dip}
\begin{enumerate}
\item[(a)] (Covenant arithmetic.) If $\mathrm{EB}_0+\Delta\mathrm{EB}_t>0$, then $\mathrm{ND}_t\le\overline{\rm lev}\,\mathrm{EB}_t$ if and only if $\Delta\mathrm{ND}_t+\overline{\rm lev}\,(-\Delta\mathrm{EB}_t)\le\mathfrak B^{\rm cov}:=(\overline{\rm lev}-{\rm lev}_0)\mathrm{EB}_0$. A dollar of transition cost that lowers tested EBITDA uses $w_{\rm cov}\overline{\rm lev}$ dollars of headroom; a dollar that is capitalized or added back uses one.
\item[(b)] (Budgets.) Under Assumption~\ref{ass:reg-dip} both budgets take the form $\widehat D(n_k)-\hat sN_{k-1}\le\mathfrak B$, with $\widehat D$ convex, strictly increasing and $\widehat D(0)=0$: public $(\widehat D,\hat s,\mathfrak B)=(D,s,\mathfrak gS)$; sponsor $\widehat D(n)=\mathcal T_1n+w_{\rm cov}\overline{\rm lev}\,D(n)$, $\hat s=w_{\rm cov}\overline{\rm lev}\,s-\mathcal T_1$, $\mathfrak B=\mathfrak B^{\rm cov}$. Let $\bar n(x):=\widehat D^{-1}(\mathfrak B+\hat sx)$ be the largest wave allowed after $x$ migrations.
\begin{enumerate}
\item[(i)] (Fastest ramp.) Let $N^g_0=0$ and $N^g_k=\min\{N,N^g_{k-1}+\bar n(N^g_{k-1})\}$ (with $\bar n$ rounded down for integer waves). Every feasible plan has $N_k\le N^g_k$ for all $k$, so it needs at least $\hat k:=\min\{k:N^g_k=N\}$ waves.
\item[(ii)] (Linear dual running.) If $D$ is linear, so that $\widehat D(n)=\hat dn$, and $\hat s>0$, then $N^g_k=\min\{N,(\mathfrak B/\hat s)[(1+\hat s/\hat d)^k-1]\}$, waves grow geometrically, $\hat k=\lceil\ln(1+\hat sN/\mathfrak B)/\ln(1+\hat s/\hat d)\rceil$, and the fastest ramp minimizes the program cost \eqref{eq:ext-progcost}.
\item[(iii)] (Convex dual running.) If $D(n)=d_1n+d_2n^2$, the unconstrained plan's dip falls wave by wave, so a violated budget binds first in wave 1. In the constrained optimum with $n_w$ waves, consecutive slack waves keep the decrement $s/(2d_2)$, and a binding wave $k$ with multiplier $\mu_k$ followed by a slack one has decrement $[s-\mu_k\widehat D'(n_k)/T_{\rm par}]/(2d_2)$, so the plan rises when $\mu_k>T_{\rm par}s/\widehat D'(n_k)$.
\end{enumerate}
\item[(c)] (Myopia as a budget.) In the continuous relaxation with $n_w$ waves, let the hard-budget problem satisfy Slater's condition with multipliers $\mu^*_k$. Every hard-budget optimum solves the soft-budget problem if $\mathfrak w T_{\rm par}\ge\max_k\mu^*_k$; if the hard-budget optimum and its multipliers are unique and $\mathfrak w T_{\rm par}<\max_k\mu^*_k$, the soft-budget optimum exceeds the guided dip in some wave. For any set of plans, the cost of the chosen plan is nondecreasing and its excess dip nonincreasing in $\mathfrak w$.
\item[(d)] (The sponsor's squeeze.) With the deadline $\bar n_w$ of Proposition~\ref{prop:reg-speed}(c), the program can be completed within covenant if and only if $N^g_{\bar n_w}=N$. Otherwise $N-N^g_{\bar n_w}$ processes miss the exit, each losing terminal value $h_xM_{\rm exit}s^{\rm E}$ ($s^{\rm E}$ its run-rate EBITDA saving), and an equity cure or covenant reset is worth up to that amount plus the cost savings it allows.
\end{enumerate}
\end{proposition}

\noindent\emph{Intuition.} A dip budget caps each wave's dip, and each completed wave loosens the cap because its savings offset the next wave's dual running. The program bootstraps: it starts small and grows, the opposite of the front-loading of Proposition~\ref{prop:ext-waves}; with linear dual running equal to the saving, waves double. Covenants price EBITDA dips at the covenant multiple, so a sponsor prefers costs that can be capitalized or added back, and its budget loosens only if a process's EBITDA at the covenant multiple exceeds its one-time cost ($\hat s>0$). Earnings pressure acts as a hard budget once its weight exceeds the budget's shadow price (the exact-penalty property of convex programs) and as a partial brake below it; surveyed executives say they would forgo value to meet earnings targets \citep{graham2005economic}.

\noindent\emph{Example (illustrative).} Take the twelve-process program of Proposition~\ref{prop:ext-waves} on \$500 million of revenue ($s=1$, $T_{\rm par}=0.5$, $d_1=0.2$, $d_2=0.25$, $O_w=0.5$, in \$ millions). Unconstrained, waves of 6, 4 and 2 move the margin by $-2.04$, $+0.24$ and $+1.72$ points at a cost of 19.70. A guided dip of one point ($\mathfrak B=\$5$ million a year) gives waves of 4, 5 and 3, with margin changes of $-0.96$, $-0.65$ and $+1.23$ points, at a cost of 20.45; this is also the fastest ramp. In the continuous relaxation the first wave's budget binds with multiplier 0.319, so the soft budget acts as a hard one for $\mathfrak w\ge0.64$, above the assumed preset; in integer plans, $\mathfrak w$ of 0 and 0.15 keep 6, 4, 2, $\mathfrak w=0.3$ chooses 5, 4, 3, and the budget plan is chosen from $\mathfrak w=0.44$. For a sponsor with EBITDA of \$100 million, entry leverage 6.1 (the top of the range) and a covenant of 6.5, headroom is \$40 million; with $\mathcal T_1=\$2.5$ million and $w_{\rm cov}=1$, $\widehat D(n)=2.5n+6.5D(n)$ and $\hat s=4$, and the optimal plan is the rising ramp 3, 4, 5 (cost 21.95). With the omitted interim cash flows it uses at most \$34.1 million and stays optimal under the exact test. If only two waves fit before the exit, 7 processes migrate and 5 miss, \$70 million of exit value at $M_{\rm exit}=14$ and $h_x=1$ (\$35 million at $h_x=0.5$); a \$10 million cure allows 4, 5, 3, and completing all twelve in two waves needs \$41.3 million (two waves of 6). With linear dual running equal to the saving and a \$5 million budget, 40 processes go in waves of 5, 10, 20 and 5.

\subsection{Impact metrics}\label{sec:regimes-metrics}

\begin{lemma}[Distribution waterfall]\label{lem:reg-waterfall}
Let a fund return the gross multiple $\mathrm{MOIC}$ on invested capital after $T$ years, with a preferred return $r_{\rm pref}$ compounded annually, carry $c_{\rm cy}$ and a full catch-up, and write $\Lambda_T:=(1+r_{\rm pref})^T$. The investors' (LPs') distributions per dollar invested are $\mathrm{MOIC}$ if $\mathrm{MOIC}\le\Lambda_T$; $\Lambda_T$ if $\Lambda_T<\mathrm{MOIC}\le1+(\Lambda_T-1)/(1-c_{\rm cy})$; and $1+(1-c_{\rm cy})(\mathrm{MOIC}-1)$ otherwise. With management fees $c_{\rm fee}$ a year on invested capital, the net multiple is these distributions divided by $1+c_{\rm fee}T$, and the net IRR of a single entry and exit is its $T$-th root less one (fees treated as paid at entry).
\end{lemma}

\begin{proposition}[Impact metrics by regime]\label{prop:reg-metrics}
\begin{enumerate}
\item[(a)] (Sponsor.) Entry equity is $\mathrm{Eq}_0=(1+\mathfrak p+c_{\rm tx})P_0-\mathrm{ND}_0$ and exit equity $\mathrm{Eq}_T=M_{\rm exit}[\mathrm{EB}^{\rm sq}_T+h_x\Delta\mathrm{EB}_T]-\mathrm{ND}_T$, where $\mathrm{EB}^{\rm sq}_T$ is EBITDA without the program, with $\mathrm{ND}_T=\mathrm{ND}_0e^{r_{\rm debt}T}-\int_0^Te^{r_{\rm debt}(T-t)}\mathrm{FCF}(t)\,dt$ when free cash flow is swept to debt. A program that adds seasoned run-rate $\Delta\mathrm{EB}_T$ and cash flows $\Delta\mathrm{CF}$ raises the gross multiple by $[h_xM_{\rm exit}\Delta\mathrm{EB}_T+\int_0^Te^{r_{\rm debt}(T-t)}d\,\Delta\mathrm{CF}(t)]/\mathrm{Eq}_0$; the gross IRR is $\mathrm{MOIC}^{1/T}-1$, and net figures follow from Lemma~\ref{lem:reg-waterfall}. Each percent of entry EBITDA added to seasoned run-rate raises the multiple by $0.01h_xM_{\rm exit}\mathrm{EB}_0/\mathrm{Eq}_0$.
\item[(b)] (Public.) The EPS path is $\Delta\mathrm{EPS}_y=(1-\tau)\Delta\mathrm{OI}_y/N_{\rm sh}$, with $\Delta\mathrm{OI}_y$ the GAAP path of Propositions~\ref{prop:tasks-scorepath}(b) and \ref{prop:ext-waves}(f); the gap to guidance is $\Delta\mathrm{OI}_y+\mathfrak gS$, and the program is accretive from the first year with $\Delta\mathrm{OI}_y\ge0$. The cash and GAAP Rule-of-40 paths are those of Proposition~\ref{prop:tasks-scorepath}.
\item[(c)] (Every regime.) To first order, $\Delta R^*_{X_{\mathsf o}}=-(1-\varphi)(\Delta c_h+\Delta c_c)-\Delta a-(g^*(X_{\mathsf o})+\delta_K)\Delta k+\sigma^*(X_{\mathsf o})\delta_\sigma+(X_{\mathsf o}-k)\Delta G_i$ (Proposition~\ref{prop:tasks-aggregate}(b) at $X=X_{\mathsf o}$; when every rival makes the savings, $-(1-\varphi)(\Delta c_h+\Delta c_c+\Delta a)$ with $\varphi=\bar\varphi^{\rm ind}C_0$ replaces the first two terms), and $\Delta\text{value}_{\mathsf o}=V_{\mathsf o}$ of the program by \eqref{eq:reg-value}; for a going concern at small share this is $S\,\Delta R^*_{X_{\mathsf o}}/(r_{\mathsf o}-g^\circ_{\mathsf o})$ less the present value of one-time costs (Proposition~\ref{prop:tasks-aggregate}(c) at $r=r_{\mathsf o}$). The regime enters $\Delta R^*$ only through the capital charge, the growth-engine term and the retention term.
\item[(d)] (Stakeholders.) At given output, a saving $W_{\rm cogs}$ in cost of revenue and $W_{\rm op}$ elsewhere gives $\bar\varphi W_{\rm cogs}$ to customers, $\tau[(1-\bar\varphi)W_{\rm cogs}+W_{\rm op}]$ in taxes and the rest to owners; displaced employees receive $(1-u_d)T_{\rm sev}C^{\rm disp}$ once.
\end{enumerate}
\end{proposition}

\noindent\emph{Intuition and example (illustrative).} The sponsor's metrics are levered: a dollar of new run-rate EBITDA is worth $h_xM_{\rm exit}$ of enterprise value at exit, so the program's uplift is affine in $h_xM_{\rm exit}$. Take a software firm with \$100 million of EBITDA growing at 5\% and free cash flow of 55\% of EBITDA, and a program that adds \$30 million of run-rate EBITDA over two years for \$40 million of one-time costs (tax rate 21\%). A sponsor buys at 14 times EBITDA and sells at 14 times after six years, so the IRR includes no multiple expansion; it pays 2\% in costs, borrows 5.2 times EBITDA at 8\% a year (7.7\% continuous) and sweeps cash to debt. With $h_x=1$, the gross multiple is 2.47 (an annual IRR of 16.2\%) against 1.71 (9.3\%) without the program, and after an 8\% preferred return, 20\% carry with catch-up and 2\% fees the investors' multiple is 1.94 (11.7\%) against 1.42 (6.0\%). With $h_x=0.5$ the gross multiple is 2.16 (13.7\%) and the investors' 1.72 (9.4\%). Each percent of entry EBITDA capitalized at exit adds 0.015 to the multiple, and the uplift is 0.76 at $h_x=1$ and 0.45 at $h_x=0.5$, of which 0.14 is interim cash; all of this is conditional on the assumed $h_x$, $T_{\rm sea}$, debt rate and costs and on a small-sample practitioner multiple. For a public firm with 100 million shares, the twelve-process program changes EPS by $-\$0.036$, $+\$0.081$ and $+\$0.095$ in its first three years in waves of 6, 4, 2, and by $-\$0.032$, $+\$0.072$ and $+\$0.095$ in the budget plan 4, 5, 3.

\subsection{When a change of ownership pays}\label{sec:regimes-lbo}

\begin{assumption}[Take-private]\label{ass:reg-lbo}
A public firm without net debt is priced at $P_0=V^{\rm sq}+V^{\rm prog}(\Pi^{\rm pub})$: the value at $r_w$ of the status quo and of the program as the public owner would execute it, plan $\Pi^{\rm pub}$, under its earnings pressure. A sponsor pays $(1+\mathfrak p)P_0$ plus costs $C_{\rm tx}=c_{\rm tx}P_0$, borrows $\mathrm{ND}_0={\rm lev}_0\mathrm{EB}_0$ at $r_{\rm debt}$, executes plan $\Pi^{\rm pe}$ with the firm's own transformation capability under its covenant and deadline, sweeps free cash flow to debt, and sells at $T$ by the EBITDA rule. With probability $p_b$ a breach or distress costs $L_b$ at $T$. The control premium required to buy the firm is $\mathfrak p^{\rm req}$.
\end{assumption}

\begin{proposition}[When a change of ownership pays]\label{prop:reg-lbo}
Under Assumption~\ref{ass:reg-lbo}, in which the buyer transforms the firm with the target's own capability (its success, cost and saving parameters are the same under either owner; Proposition~\ref{prop:acq-owner} relaxes this, and an owner with its own capability adds $\Delta V^{\rm cap}_{\mathsf o}/P_0$ to $\mathfrak p^{\max}$, \eqref{eq:acq-pmax}):
\begin{enumerate}
\item[(a)] (Sponsor's value.) The sponsor's NPV at its hurdle is
\begin{equation}\label{eq:reg-lbo}
\mathrm{NPV}^{\rm sp}=V^{\rm sp}(\Pi^{\rm pe})-(1+\mathfrak p)P_0-C_{\rm tx}+\mathfrak L-e^{-r_{\rm pe}T}p_bL_b,\qquad \mathfrak L:=\mathrm{ND}_0\bigl[1-e^{-(r_{\rm pe}-r_{\rm debt})T}\bigr],
\end{equation}
where $V^{\rm sp}(\Pi):=e^{-r_{\rm pe}T}[M_{\rm exit}(\mathrm{EB}^{\rm sq}_T+h_x\Delta\mathrm{EB}_T(\Pi))+\int_0^Te^{r_{\rm debt}(T-t)}\mathrm{FCF}(t;\Pi)\,dt]$. The take-private pays if and only if $\mathfrak p^{\max}:=\{V^{\rm sp}(\Pi^{\rm pe})+\mathfrak L-C_{\rm tx}-e^{-r_{\rm pe}T}p_bL_b\}/P_0-1\ge\mathfrak p^{\rm req}$; since $\mathrm{NPV}^{\rm sp}$ is decreasing in $r_{\rm pe}$ when exit equity is positive, the break-even hurdle is the deal's continuously compounded IRR.
\item[(b)] (Frictionless benchmark.) If $r_{\rm pe}=r_{\rm debt}=r_w$ and the exit price equals the value at $T$, at $r_w$, of the firm's remaining cash flows under $\Pi^{\rm pe}$, then
\begin{equation}\label{eq:reg-exec}
\mathfrak p^{\max}P_0=\Delta V^{\rm exec}-C_{\rm tx}-e^{-r_wT}p_bL_b,\qquad \Delta V^{\rm exec}:=V^{\rm prog}(\Pi^{\rm pe})-V^{\rm prog}(\Pi^{\rm pub}),
\end{equation}
both at $r_w$: the change of ownership pays for the premium if and only if executing the transformation privately adds more value than the premium, the deal costs and the expected distress costs. The premium and $\mathfrak L$ are transfers among claimants.
\item[(c)] (Sources of the execution gain.) $\Delta V^{\rm exec}$ is nondecreasing in the public owner's myopia weight $\mathfrak w$, because the market value of the plan the public owner chooses is nonincreasing in $\mathfrak w$. It falls with the sponsor's covenant and deadline (Proposition~\ref{prop:reg-dip}(d)) and with the sponsor's lower growth weight: replacing $X$ by $X^{\rm pe}$ costs $\Pi(X)[1-(\rho-\alpha\rho^{1/\alpha})/(1-\alpha)]$ of the true criterion a year, with $\rho=[(X^{\rm pe}-k)/(X-k)]^{\alpha/(1-\alpha)}$.
\item[(d)] (Away from the benchmark.) In general $\mathfrak p^{\max}P_0=\Delta V^{\rm exec}+\mathfrak D+\mathfrak L-C_{\rm tx}-e^{-r_{\rm pe}T}p_bL_b$, where the valuation wedge $\mathfrak D:=V^{\rm sp}(\Pi^{\rm pe})-[V^{\rm sq}+V^{\rm prog}(\Pi^{\rm pe})]$ compares the sponsor's valuation, at its hurdle and exit multiple, with the value at $r_w$ of the firm it will run.
\end{enumerate}
\end{proposition}

\noindent\emph{Intuition.} Without differences in rates and exit prices, a take-private is a bet that private ownership executes better, and the premium must come out of that difference. Executing the same program a few years faster is worth little; avoiding a public owner's refusal to bear the dip is worth the whole program. A hurdle above the market rate lowers what a sponsor can pay, and debt cheaper than the hurdle raises it, a transfer rather than value created. Going-private plants restructure faster \citep{bharath2014do}, experienced sponsors improve operations \citep{bernstein2016operational}, and effects differ by deal type and sponsor, with employment falling after public-to-private buyouts and rising after private-to-private ones \citep{davis2019heterogenous}, so presets are ranges; public firms' lower, less responsive investment \citep{asker2015corporate} fits the brake of part (c).

\noindent\emph{Example (illustrative).} Let the firm of the previous example be public, worth \$1.40 billion at $r_w$ (14.0 times EBITDA) without the program, which is worth \$551 million at $r_w$ over two years and \$522 million over five. If the public owner merely executes it more slowly, $\Delta V^{\rm exec}=\$29$ million and the frictionless maximum premium is $-0.5\%$; if earnings pressure blocks it (price \$1.40 billion), $\Delta V^{\rm exec}=\$551$ million and the maximum premium is 37\%, above the required 30\%. At the sponsor's terms, with the exit multiple equal to the 18.2 times it pays, a 19.9\% hurdle (22\% annual) and 7.7\% debt, the leverage term is $\mathfrak L=\$270$ million, the valuation wedge is $\mathfrak D=-\$798$ million, and $\mathfrak p^{\max}=-0.4\%$ with $h_x=1$ ($-8.3\%$ with $h_x=0.5$). The deal pays at a 30\% premium only for a hurdle up to 13.5\% (14.4\% annual; 11.3\% and 11.9\% with $h_x=0.5$), below the 22\% survey target, conditional on the assumed $h_x$, debt rate and costs. This firm has no growth engine; for the firm of Proposition~\ref{prop:reg-weight}, moving its weight from the model's 8 to the sponsor's 2.70 would cost 2.41 points of its true criterion a year, and from illustrative public weights of \regXbOne, \regXbTwo{} and \regXbThree, \regXbLossPEOne, \regXbLossPETwo{} and \regXbLossPEThree{} points; at \regXbOne{} the sponsor's weight from mid-hold on (2.70) is the higher one, so from then a take-private would raise growth spend, not cut it. % XBOUNDS

\subsection{Hybrid regimes}\label{sec:regimes-hybrid}

\begin{proposition}[Hybrids]\label{prop:reg-hybrid}
\begin{enumerate}
\item[(a)] (Controlling holder or dual-class shares.) A public firm with myopia weight $\mathfrak w_{\rm dc}\in[0,\mathfrak w_{\rm pub}]$ has growth weight $X^{\rm my}(\mathfrak w_{\rm dc})$ between $X^{\rm my}(\mathfrak w_{\rm pub})$ and $X$, and the cost and dip of its program lie between those of the widely held public firm and the private firm (Proposition~\ref{prop:reg-dip}(c)).
\item[(b)] (Minority stake or PIPE.) (i) New equity $C_P$ raises a covenant or cash budget by $C_P$ and leaves a GAAP budget unchanged, so it speeds up a financing-constrained program and not a pressure-constrained one. (ii) Because $R_X$ is linear in $X$, a board that weights an investor's per-date criterion by $\omega_m$ and the public holders' by $1-\omega_m$ acts on the weight $\omega_mX^{\rm inv}+(1-\omega_m)X^{\rm pub}$.
\item[(c)] (Carve-out.) If the separated unit must otherwise rebuild a shared process at $C^{\rm sep}$, with a parallel run under a transition services agreement costing $C^{\rm par,sh}$ a year for $T_{\rm par}$, automation's incremental value is $\mathrm{NPV}_p+C^{\rm sep}+C^{\rm par,sh}\mathcal A_r(T_{\rm par})$, and without discounting the minimum horizon of Proposition~\ref{prop:tasks-transition}(b) falls by $(C^{\rm sep}+C^{\rm par,sh}T_{\rm par})/(\varpi_p\bar b_pW)$.
\item[(d)] (Venture-backed growth firm.) The weight is \eqref{eq:reg-Xexit} with $r_{\rm vc}$ and $X_T=(V/S)_T$; at exit, by Lemma~\ref{lem:valuation2}, the value changes by $\mathfrak A_TS_T(\Delta m_T+X^{\rm scr}\Delta g_T)$ to first order, so the terminal objective is the exit-date Rule-of-$X^{\rm scr}$ score.
\end{enumerate}
\end{proposition}

\noindent\emph{Intuition and example (illustrative).} Hybrids mix the pure regimes' parameters, and the results are monotone in each, so a hybrid's plan lies between its components'. Financing and disclosure act separately: a PIPE relaxes a cash constraint, not one on reported earnings. A carve-out changes the counterfactual: when a process must be rebuilt anyway, the rebuild and the parallel run are shared by both routes. With the assumed $\mathfrak w_{\rm dc}=0.15$, the dual-class firm's weight is 6.96 against 6.15 widely held and 8 private, and it keeps the private plan 6, 4, 2 where the widely held firm chooses 5, 4, 3. A \$10 million equity injection leaves the public budget plan at 4, 5, 3 and moves the covenant plan from 3, 4, 5 to 4, 5, 3. If a carve-out must otherwise rebuild Case R's process at \$600{,}000, with its \$800{,}000-a-year duplication premium borne on both routes, automating gains \$735{,}000 and its minimum horizon falls from 9.9 to 3.7 years. A venture-backed firm seven years from exit whose marginal revenue earns 15\%, grows 25\% and sells at 8 times revenue has $X^{\rm x}=7.93$ at 27.0\% (31\% annual), a target rate, since venture returns depend on correcting for selected valuations \citep{cochrane2005risk,korteweg2010risk}.
